\documentclass[11pt, oneside]{article} 
\usepackage{geometry}
\geometry{letterpaper} 
\usepackage{graphicx}
	
\usepackage{amssymb}
\usepackage{amsmath}
\usepackage{parskip}
\usepackage{color}
\usepackage{hyperref}

\graphicspath{{/Users/telliott/Dropbox/Github-math/figures/}}
% \begin{center} \includegraphics [scale=0.4] {gauss3.png} \end{center}

\title{Gardner's angles}
\date{}

\begin{document}
\maketitle
\Large

%[my-super-duper-separator]

Here is a puzzle I came across on the internet.

\url{https://www.theguardian.com/science/alexs-adventures-in-numberland/2014/oct/21/martin-gardner-mathematical-puzzles-birthday}

It is attributed to Martin Gardner.

\begin{center} \includegraphics [scale=0.35] {gardner1.png} \end{center}

\emph{To prove}.

Use elementary geometry to show that $\angle A + \angle B = \angle C$.
\begin{center} \includegraphics [scale=0.35] {gardner2.png} \end{center}

I see two equations to write immediately:
\[ B = A + x \]
\[ C = A + x + y \]
\[ = B + y \]

If it is really true that $C = A + B$, then $y$ must be equal to $A$, and also, $x + y$ must be equal to $B$.  

Alas, I did not see a way to demonstrate that using algebra alone.  I had the solution in front of me, temptation got the better of me, and I peeked.

\subsection*{trigonometry}

To make up for this lapse, let's first just use some trigonometric tools to explore the problem.  Let each square be a unit square.  Clearly $\tan C = 1, \tan B = 1/2$ and $\tan A = 1/3$.

Recall the sum of angles formulas:
\[ \sin s + t = \sin s \cos t + \cos s \sin t \]
\[ \cos s + t = \cos s \cos t - \sin s \sin t \]

So
\[ \tan s + t = \frac{\sin s \cos t + \cos s \sin t}{\cos s \cos t - \sin s \sin t} \]
Multiply through by $1/\cos s \cos t$ on top and bottom to obtain
\[ \tan s + t = \frac{\tan s + \tan t}{1 - \tan s \tan t} \]

\[ \tan A = \frac{1/2 + 1/3}{1 - 1/2 \cdot 1/3} \]
which does indeed equal $1$.  

Since the inverse tangent is a single-valued function over the interval $[0,\pi/2]$, and $\tan A = \tan B + C$, it follows that $\angle A = \angle B + \angle C$.

\subsection*{Pythagorean theorem}

We can also get some help from Pythagoras.  Divide the unit square into four small squares with side lengths equal to one-half.  Draw the diagonals in one of those smaller squares.

\begin{center} \includegraphics [scale=0.4] {gardner5.png} \end{center}

The resulting small triangles, formed by crossing diagonals in a square, are all congruent.  We don't need to do any algebra.

Thus, $a$ is one-quarter of the length of the diagonal.

The angle which we called $y$ before, and is here the upper of the two angles marked with black dots, lies in a right triangle with opposite side $a$ and adjacent side $3a$, so its tangent is $1/3$.  

This is the same as the tangent of the angle we called $A$ before, which here is the lower of the two marked angles.  That is just as we expect.  We have established that the two marked angles are part of similar right triangles, and therefore equal to each other.

\subsection*{geometric solution}

The geometric solution is to draw another rectangle.
\begin{center} \includegraphics [scale=0.35] {gardner3.png} \end{center}

The long side of the rectangle is equal to $2 \sqrt{2}$ and is twice the short side, which is just the diagonal of the original unit square.

Therefore, the right triangle containing $\angle D$ is similar to the right triangle containing $\angle B$.  We conclude that $\angle B = \angle D$.

Since $\angle A + \angle D = \angle C$, it follows that $\angle A + \angle B = \angle C$.

$\square$

\end{document}